\documentclass[10pt,a4paper]{amsart}
\usepackage{amsmath,amssymb,amsthm,graphicx}
% Retained minimal preamble for the Stacks Project chapter more-algebra.tex (More on Algebra).
% The chapter file pinned-original-source.tex carries only \input{preamble} on its first line, so no
% preamble text lives inside the chapter itself.  The Stacks Project uses one single shared preamble for
% every chapter; the snapshot stages it identically as sources/stacks-chow/preamble.tex and
% sources/stacks-etale-cohomology/preamble.tex.  The block below is copied mechanically, line for line,
% from sources/stacks-chow/preamble.tex lines 14-18 (the Xy-pic setup) and 154-217 (the native theorem
% declarations together with the mathematical macros).  Nothing was added, removed or redefined; the
% wrapper class and the amsmath/amssymb/amsthm/graphicx base packages are supplied by the assembler.
% Chapter: more-algebra.tex -- More on Algebra.
% For commutative diagrams we use Xy-pic
\usepackage[all]{xy}

% We use 2cell for 2-commutative diagrams.
\xyoption{2cell}

% Theorem environments.
%
\theoremstyle{plain}
\newtheorem{theorem}[subsection]{Theorem}
\newtheorem{proposition}[subsection]{Proposition}
\newtheorem{lemma}[subsection]{Lemma}

\theoremstyle{definition}
\newtheorem{definition}[subsection]{Definition}
\newtheorem{example}[subsection]{Example}
\newtheorem{exercise}[subsection]{Exercise}
\newtheorem{situation}[subsection]{Situation}

\theoremstyle{remark}
\newtheorem{remark}[subsection]{Remark}
\newtheorem{remarks}[subsection]{Remarks}

\numberwithin{equation}{subsection}

% Macros
%
\def\lim{\mathop{\mathrm{lim}}\nolimits}
\def\colim{\mathop{\mathrm{colim}}\nolimits}
\def\Spec{\mathop{\mathrm{Spec}}}
\def\Hom{\mathop{\mathrm{Hom}}\nolimits}
\def\Ext{\mathop{\mathrm{Ext}}\nolimits}
\def\SheafHom{\mathop{\mathcal{H}\!\mathit{om}}\nolimits}
\def\SheafExt{\mathop{\mathcal{E}\!\mathit{xt}}\nolimits}
\def\Sch{\mathit{Sch}}
\def\Mor{\mathop{\mathrm{Mor}}\nolimits}
\def\Ob{\mathop{\mathrm{Ob}}\nolimits}
\def\Sh{\mathop{\mathit{Sh}}\nolimits}
\def\NL{\mathop{N\!L}\nolimits}
\def\CH{\mathop{\mathrm{CH}}\nolimits}
\def\proetale{{pro\text{-}\acute{e}tale}}
\def\etale{{\acute{e}tale}}
\def\QCoh{\mathit{QCoh}}
\def\Ker{\mathop{\mathrm{Ker}}}
\def\Im{\mathop{\mathrm{Im}}}
\def\Coker{\mathop{\mathrm{Coker}}}
\def\Coim{\mathop{\mathrm{Coim}}}

% Boxtimes
%
\DeclareMathSymbol{\boxtimes}{\mathbin}{AMSa}{"02}

%
% Macros for moduli stacks/spaces
%
\def\QCohstack{\mathcal{QC}\!\mathit{oh}}
\def\Cohstack{\mathcal{C}\!\mathit{oh}}
\def\Spacesstack{\mathcal{S}\!\mathit{paces}}
\def\Quotfunctor{\mathrm{Quot}}
\def\Hilbfunctor{\mathrm{Hilb}}
\def\Curvesstack{\mathcal{C}\!\mathit{urves}}
\def\Polarizedstack{\mathcal{P}\!\mathit{olarized}}
\def\Complexesstack{\mathcal{C}\!\mathit{omplexes}}
% \Pic is the operator that assigns to X its picard group, usage \Pic(X)
% \Picardstack_{X/B} denotes the Picard stack of X over B
% \Picardfunctor_{X/B} denotes the Picard functor of X over B
\def\Pic{\mathop{\mathrm{Pic}}\nolimits}
\def\Picardstack{\mathcal{P}\!\mathit{ic}}
\def\Picardfunctor{\mathrm{Pic}}
\def\Deformationcategory{\mathcal{D}\!\mathit{ef}}

% --- appended 2026-10-12 (Chief): the project own definitions of the reference/slogan/history
% display environments, copied mechanically from the shared Stacks preamble lines 8-12
% (\usepackage{verbatim} plus the three \newenvironment lines, which the chapter preamble keeps
% verbatim from the project preamble). Nothing added, removed or redefined: these environments
% carry no mathematical content and the project itself typesets them as comments. Without these
% three lines a counted statement whose own text uses them fails with Environment slogan undefined.
\usepackage{verbatim}
\newenvironment{reference}{\comment}{\endcomment}
\newenvironment{slogan}{\comment}{\endcomment}
\newenvironment{history}{\comment}{\endcomment}
\usepackage{cleveref}
\crefname{originallemma}{original lemma}{original lemmas}
\Crefname{originallemma}{Original lemma}{Original lemmas}
\makeatletter
\newlabel{lemma-composition}{{17966}{source}{}{sourcepointer-1}{}}
\newlabel{lemma-composition@cref}{{[originallemma][0][]17966}{[1][0][]source}{}{}{}}
\makeatother
\title{Composition of internal Hom admits a canonical functorial morphism in the derived category}
\author{The Stacks Project Authors}
\date{stacks-project:more-algebra.tex:local-snapshot-20261012; GNU FDL 1.2 or later}
\begin{document}
\maketitle
\setcounter{subsection}{1}
\noindent\textit{Editorial note.} All delivered mathematical text is the author's own, copied from the pinned Stacks Project chapter file more-algebra.tex; no statement or proof step was rewritten, reordered or removed.\par\medskip
\noindent\textit{Editorial note.} The chapter file pinned-original-source.tex carries only \textbackslash{}input\{preamble\} on its first line, so no preamble text lives inside the chapter. The retained file sources/stacks-more-algebra/retained-preamble.tex is a mechanical, line-for-line copy of the project's single shared preamble as staged in the snapshot; nothing was added, removed or redefined and no package was added.\par\medskip
\noindent\textit{Editorial note.} The source URL is the source-level Stacks Project URL, as required when more than one card is produced from one source; the per-claim tag is carried by claim\_key.\par\medskip
\noindent\textit{Editorial note.} The author writes 'We omit the proof that the construction is functorial.' at the end of the counted proof span; that sentence is kept verbatim and marks the step left to the reader. The single reference inside the counted proof span is recorded as a non-load-bearing navigation pointer.\par\medskip
\noindent\textit{Allowed standard prerequisites (editorial).}\begin{itemize}
\item The composition map for Hom complexes of K-injective representatives. invoked at proof lines 18771-18778. Is the underlying map of complexes whose lift to the derived category is constructed.
\item K-flat replacements and the fact that a K-flat complex computes the derived tensor product. used at proof lines 18782-18793. Is what makes the source of the first displayed arrow represent the derived tensor product.
\item Hom complexes of K-injective representatives representing the derived Hom. used at proof lines 18779-18781. Identifies the three Hom complexes with the objects named in the statement.
\end{itemize}\medskip
\ifdefined\hypertarget\hypertarget{sourcepointer-1}{}\fi
\noindent\textit{Source reference.} Original lemma 17966: sources/stacks-more-algebra/pinned-original-source.tex, source lines 17966--17966. Same-chapter navigation pointer; the statement it names is standard for the section and the argument that consumes it stays inside the counted proof span.\par\medskip
\begin{lemma}
\label{lemma-internal-hom-composition}
Let $R$ be a ring. Given $K, L, M$ in $D(R)$ there is a canonical morphism
$$
R\Hom_R(L, M) \otimes_R^\mathbf{L} R\Hom_R(K, L) \longrightarrow R\Hom_R(K, M)
$$
in $D(R)$ functorial in $K, L, M$.
\end{lemma}

\begin{proof}
Choose a K-injective complex $I^\bullet$ representing $M$,
a K-injective complex $J^\bullet$ representing $L$, and
any complex of $R$-modules $K^\bullet$ representing $K$.
By Lemma \ref{lemma-composition} there is a map of complexes
$$
\text{Tot}\left(
\Hom^\bullet(J^\bullet, I^\bullet) \otimes_R \Hom^\bullet(K^\bullet, J^\bullet)
\right)
\longrightarrow
\Hom^\bullet(K^\bullet, I^\bullet)
$$
The complexes of $R$-modules $\Hom^\bullet(J^\bullet, I^\bullet)$,
$\Hom^\bullet(K^\bullet, J^\bullet)$, and $\Hom^\bullet(K^\bullet, I^\bullet)$
represent $R\Hom_R(L, M)$, $R\Hom_R(K, L)$, and $R\Hom_R(K, M)$.
If we choose a K-flat complex $H^\bullet$ and a quasi-isomorphism
$H^\bullet \to \Hom^\bullet(K^\bullet, J^\bullet)$, then there is a map
$$
\text{Tot}\left(
\Hom^\bullet(J^\bullet, I^\bullet) \otimes_R H^\bullet
\right)
\longrightarrow
\text{Tot}\left(
\Hom^\bullet(J^\bullet, I^\bullet) \otimes_R \Hom^\bullet(K^\bullet, J^\bullet)
\right)
$$
whose source represents $R\Hom_R(L, M) \otimes_R^\mathbf{L} R\Hom_R(K, L)$.
Composing the two displayed arrows gives the desired map. We omit the
proof that the construction is functorial.
\end{proof}

\end{document}
